\section{方法与比较}
\label{sec:method}
\subsection{统一比较口径}
假设两种方案具有相同使用时间，用下式记录相对变化：
\begin{equation}
  r=\frac{E_{\mathrm{A}}-E_{\mathrm{B}}}{E_{\mathrm{A}}}\times100\%.
  \label{eq:change}
\end{equation}
\begin{table}[htbp]
\centering
\caption{虚构的教学数据，不代表真实测量结果}
\label{tab:energy}
\begin{tabular}{lrr}
\toprule
方案 & 使用时间（小时） & 能耗（千瓦时） \\
\midrule
甲 & 8 & 12.0 \\
乙 & 8 & 9.6 \\
\bottomrule
\end{tabular}
\end{table}
将表\ref{tab:energy}中的数字代入式\ref{eq:change}，得到示意变化率。
真实研究还需要注明测量条件、仪器和误差。
